\Needspace{9\baselineskip}\begin{center}\textit{Conversion of common fractions into decimals.}\end{center}

\Needspace{7\baselineskip}
\Article{312}
\textsc{Definition.} When a common fraction is converted into a decimal, we call the sequence of decimal digits\HistoricalNote{*)}{For brevity we restrict the following investigation to the usual decimal system, since it can easily be extended to any other.} (excluding the integer part, if there is one), whether finite or continuing indefinitely, the fraction's \textit{mantissa}. We use this expression, otherwise customary only for logarithms, in a wider sense. Thus, \textit{for example}, the mantissa of the fraction \(\frac18\) is \(125\), that of \(\frac{35}{16}\) is \(1875\), and that of \(\frac2{37}\) is \(054054\ldots\), continuing indefinitely.

It follows at once from this definition that the fractions \(\frac ln,\frac mn\), which have the same denominator, have the same or different mantissas according as their numerators \(l,m\) are congruent or incongruent modulo \(n\). A finite mantissa does not change if any number of zeros are added on the right. The mantissa of the fraction \(\frac{10m}{n}\) is obtained by removing the first digit from the mantissa of \(\frac mn\). More generally, the mantissa of \(\frac{10^\nu m}{n}\) is found by removing the first \(\nu\) digits from that of \(\frac mn\). The mantissa of \(\frac1n\) begins at once with a significant digit, \textit{that is}, one different from zero, if \(n\) is not \(>10\). But if \(n>10\), and the denominator is unequal to any power of \(10\) and has \(k\) digits, the first \(k-1\) digits of the mantissa of \(\frac1n\) are zeros, and only the next, the \(k\)-th, is significant. It follows easily that if \(\frac ln,\frac mn\) have different mantissas, \textit{that is},
\SourceLeaf{388}
if \(l,m\) are incongruent modulo \(n\), these mantissas certainly cannot agree in their first \(k\) digits, but must differ at or before the \(k\)-th digit.

\Needspace{7\baselineskip}
\Article{313}
\textsc{Problem.} \textit{Given the denominator of a fraction \(\frac mn\) and the first \(k\) digits of its mantissa, find the numerator \(m\), which we suppose less than \(n\).}

\textit{Solution.} Regard those \(k\) digits as an integer, multiply it by \(n\), and divide the product by \(.10^k\), or remove the last \(k\) digits. If the quotient is an integer, that is, if the digits removed are zeros, it is clearly the required numerator and the given mantissa is complete. Otherwise the required numerator is the next greater integer: the quotient increased by one after the following decimal digits have been discarded. The reason for this rule is so easily seen from the observations at the end of the preceding article that a fuller explanation is unnecessary.

\textit{Example.} If the first two digits of the mantissa of a fraction with denominator \(23\) are known to be \(69\), we have the product \(23.69=1587\). Removing its last two digits and adding one gives the required numerator \(=16\).

\Needspace{7\baselineskip}
\Article{314}
We begin with fractions whose denominators are primes or powers of primes, and will later show how the others can be reduced to these. First we observe at once that the mantissa of a fraction \(\frac a{p^\mu}\), whose numerator \(a\) we always suppose not divisible by the prime \(p\), is finite and has \(\mu\) digits if \(p=2\) or \(=5\). In the former case this mantissa, regarded as an integer, \(=5^\mu a\), and in the latter, \(=2^\mu a\). These facts are so obvious that they need no explanation.

If \(p\) is any other prime, however, \(10^r a\) is never divisible by \(p^\mu\), however large \(r\) is taken. Thus the mantissa of \(F=\frac a{p^\mu}\) must continue indefinitely. Suppose \(10^e\) is the lowest power of \(10\) congruent to unity modulo \(p^\mu\). Compare Section III, where we showed that \(e\) either equals \((p-1)p^{\mu-1}\) or is a proper divisor of it. It is easy to see that \(10^e a\) is also the first number in the series \(10a,100a\), \(1000a\), and so on, congruent to \(a\) modulo the same number. Now, by Article \(312\), the mantissas of the fractions \(\frac{10a}{p^\mu},\frac{100a}{p^\mu}\ldots\frac{10^e a}{p^\mu}\) arise by removing from the mantissa
\SourceLeaf{389}
of \(F\) the first digit, the first two \(\ldots e\) digits, respectively. So the same digits clearly repeat in this mantissa after the first \(e\) digits, and not sooner. We may call these first \(e\) digits, whose infinite repetition forms the mantissa, the \textit{period} of this mantissa or of the fraction \(F\). Clearly the length of the period, \textit{that is}, the number of digits it contains (\(=e\)), is entirely independent of the numerator \(a\), and is determined by the denominator alone. Thus, \textit{for example}, the period of \(\frac1{11}\) is \(09\), and that of \(\frac37\) is \(428571\)\HistoricalNote{*)}{The distinguished Robertson marks the beginning and end of the period by placing two dots above its first and last digits (\textit{Theory of circulating fractions, Philos. Trans.} \(1769\) \textit{p.} \(207\)), which we do not think necessary here.}.

\Needspace{7\baselineskip}
\Article{315}
Once the period of a fraction is known, its mantissa can therefore be extended to any number of digits. Further, if \(b\equiv10^\lambda a\; (\mathrm{mod}.\ p^\mu)\), the period of \(\frac b{p^\mu}\) clearly results from placing the first \(\lambda\) digits of the period of \(F\), assuming \(\lambda<e\), as we may, after the remaining \(e-\lambda\) digits. Thus, with the period of \(F\) we also have the periods of all fractions whose numerators are congruent to \(10a,100a,1000a\), and so on, modulo the denominator \(p^\mu\). Thus, \textit{for example}, since \(6\equiv3.10^2\; (\mathrm{mod}.\ 7)\), the period of \(\frac67\) is obtained at once from that of \(\frac37\) as \(857142\).

Whenever \(p^\mu\) is the modulus and \(10\) is a primitive root, by Articles \(57,89\), the period of \(\frac1{p^\mu}\) immediately gives the period of any other fraction \(\frac m{p^\mu}\) whose numerator \(m\) is not divisible by \(p\). We remove as many digits from its left and restore them on its right as there are units in the index of \(m\), with \(10\) taken as base. This makes clear why, in this case, \(10\) has always been taken as base in Table I; see Article \(72\).

\begin{samepage}
But when \(10\) is not a primitive root, the period of \(\frac1{p^\mu}\) gives only the periods of fractions whose numerators are congruent to some power of \(10\) modulo \(p^\mu\). Let \(10^e\) be the lowest power of \(10\) congruent to unity modulo \(p^\mu\), let \((p-1)p^{\mu-1}=ef\), and take as base a primitive root \(r\) for which the index of \(10\) is \(f\), by Article \(71\). In this system the numerators of the fractions whose periods can be obtained from that of \(\frac1{p^\mu}\) have indices \(f,2f,3f\ldots ef-f\). Similarly, the period of \(\frac r{p^\mu}\) gives the periods of fractions whose numerators \(10r,100r,1000r\), and so on, correspond to the indices \(f+1,2f+1,3f+1\), and so on. The period of the fraction with numerator \(rr\), whose index is \(2\), gives the periods of fractions whose numerators
\SourceLeaf{390}
have indices \(f+2,2f+2,3f+2\), and so on. More generally, the period of the fraction with numerator \(r^i\) gives the periods of fractions whose numerators have indices \(f+i,2f+i,3f+i\), and so on. It follows easily that, once we have the periods of the fractions with numerators \(1,r,rr,r^3\ldots r^{f-1}\), all the others can be deduced by transposition alone, using the following rule. Let the index of the numerator \(m\) of the given fraction \(\frac m{p^\mu}\), in the system with base \(r\), \(=i\), which we suppose less than \((p-1)p^{\mu-1}\). By division by \(f\), put \(i=\alpha f+\beta\), with \(\alpha,\beta\) positive integers, or also \(0\), and \(\beta<f\). Then the period of \(\frac m{p^\mu}\) is obtained from the period of the fraction with numerator \(r^\beta\), and so \(1\) when \(\beta=0\), by placing its first \(\alpha\) digits after the others. Thus we retain that very period when \(\alpha=0\). This sufficiently explains why, in constructing Table I, we followed the rule explained in Article \(72\).
\par\end{samepage}
